\scanpage{289}
On considère les foncteurs d'inclusion
\[
(9)\qquad i_N:\mathcal M_N\hookrightarrow\mathcal M\qquad
 i^N_{\deleted{\unsure{N}}}:\mathcal M^N\unsure{-|\longrightarrow}\mathcal M
\]
plus généralement, pour $N\leq N'$
\[
(9.1)\qquad i_{N,N'}:\mathcal M_{N'}\hookrightarrow\mathcal M_N\qquad
 i^{N',N}:\mathcal M^N\subset\mathcal M^{N'}
\]
\deleted{On va étudier les}

\unsure{donc} :
\[
(9.2)\qquad i_Ni_{N,N'}=i_{N'\deleted{\unsure{,N}}}\quad,\qquad
 i^{N'}i^{N',N}=i^N
\]
\deleted{\unsure{plus g.}} et si $N\leq N'\leq N''$,\quad \unsure{d'où} \deleted{$\unsure{i_N}$}
\[
(9.3)\quad\left\{\begin{array}{ll}
\mathcal M_{N''}\subset\mathcal M_{N'}\subset\mathcal M_N\ \deleted{\text{et}}&\mathcal M^N\subset\mathcal M^{N'}\subset\mathcal M^{N''}\\[1ex]
i_{N,N''}=i_{N,N'}i_{N',N''}&i^{N'',N}=i^{N'',N'}i^{N',N}
\end{array}\right.
\]
On va étudier la suite de foncteurs adjoints associés à $i_N,i_{N,N'}$ d'une part, $i^N,i^{N,N'}$ d'autre part.

\uline{Proposition} 1 $i_N$\inserted{$:\mathcal M_N\hookrightarrow\mathcal M$} s'insère dans une suite de 4 foncteurs adjoints
\[
(10)\qquad q_N\quad i_N\quad\tau_N\quad\varphi_N
\]
(chaque foncteur \deleted{\unsure{est}} parmi $i_N,\tau_N,\varphi_N$ est adjoint à droite de celui écrit à sa gauche)

définis ainsi :
\[
a)\quad\tau_N:\mathcal M\longrightarrow\mathcal M_N\qquad
\text{(foncteur tronquage à droite)}
\]
\[
\tau_N(P)\inserted{i_N\tau_N(P)}\subset\deleted{\unsure{P}}P\qquad\text{(sous objet de }P\text{)}.
\]
\margin{hom d'adjonction}
\[
(10.1)\qquad(\tau_N(P))_i=\left\{\begin{array}{ll}P_i&\text{si }i\geq N\\0&\text{si }i<N\end{array}\right.
\]
\margin{290}
